\begin{table*}[t]\centering\caption{Main results: test IoU (mean $\pm$ std over seeds) at $N=500$. Bold: best, underline: second-best per task. Band-ignore CE at flip $p=0.1,0.2$ has three seeds; there is no CE/GCE/SCE row for those tasks in this group (CE is in Table~\ref{tab:noise}).}\label{tab:main}
\begin{tabular}{llccccc}
\toprule
Method & Task & miou $\uparrow$ & iou\_circle $\uparrow$ & iou\_square $\uparrow$ & iou\_triangle $\uparrow$ & $n$ \\
\midrule
CE & none & \textbf{0.80 $\pm$ 0.03} & \textbf{0.82 $\pm$ 0.03} & \underline{0.71 $\pm$ 0.04} & \textbf{0.67 $\pm$ 0.07} & 4 \\
SCE & none & 0.77 $\pm$ 0.02 & 0.80 $\pm$ 0.02 & 0.69 $\pm$ 0.02 & 0.60 $\pm$ 0.08 & 4 \\
GCE & none & \underline{0.79 $\pm$ 0.02} & \underline{0.81 $\pm$ 0.01} & \textbf{0.71 $\pm$ 0.01} & \underline{0.67 $\pm$ 0.05} & 4 \\
\textbf{Band-ignore CE} & none & 0.64 $\pm$ 0.07 & 0.65 $\pm$ 0.07 & 0.47 $\pm$ 0.19 & 0.48 $\pm$ 0.05 & 4 \\
\midrule
CE & boundary\_p0.3 & \underline{0.78 $\pm$ 0.04} & \textbf{0.81 $\pm$ 0.04} & \underline{0.69 $\pm$ 0.03} & \textbf{0.65 $\pm$ 0.09} & 4 \\
SCE & boundary\_p0.3 & 0.70 $\pm$ 0.09 & 0.80 $\pm$ 0.03 & 0.65 $\pm$ 0.07 & 0.38 $\pm$ 0.27 & 4 \\
GCE & boundary\_p0.3 & \textbf{0.78 $\pm$ 0.02} & \underline{0.81 $\pm$ 0.02} & \textbf{0.70 $\pm$ 0.02} & \underline{0.64 $\pm$ 0.04} & 4 \\
\textbf{Band-ignore CE} & boundary\_p0.3 & 0.59 $\pm$ 0.08 & 0.63 $\pm$ 0.08 & 0.42 $\pm$ 0.22 & 0.36 $\pm$ 0.13 & 4 \\
\midrule
CE & flip\_p0.3 & \underline{0.53 $\pm$ 0.10} & 0.54 $\pm$ 0.09 & 0.28 $\pm$ 0.14 & \underline{0.32 $\pm$ 0.19} & 4 \\
SCE & flip\_p0.3 & 0.53 $\pm$ 0.13 & \underline{0.57 $\pm$ 0.12} & \underline{0.36 $\pm$ 0.23} & 0.21 $\pm$ 0.25 & 4 \\
GCE & flip\_p0.3 & \textbf{0.59 $\pm$ 0.09} & \textbf{0.60 $\pm$ 0.07} & \textbf{0.43 $\pm$ 0.08} & \textbf{0.38 $\pm$ 0.25} & 4 \\
\textbf{Band-ignore CE} & flip\_p0.3 & 0.45 $\pm$ 0.09 & 0.46 $\pm$ 0.06 & 0.23 $\pm$ 0.15 & 0.17 $\pm$ 0.19 & 4 \\
\midrule
\textbf{Band-ignore CE} & flip\_p0.1 & \textbf{0.52 $\pm$ 0.12} & \textbf{0.54 $\pm$ 0.10} & \textbf{0.31 $\pm$ 0.25} & \textbf{0.28 $\pm$ 0.15} & 3 \\
\midrule
\textbf{Band-ignore CE} & flip\_p0.2 & \textbf{0.44 $\pm$ 0.09} & \textbf{0.46 $\pm$ 0.06} & \textbf{0.22 $\pm$ 0.19} & \textbf{0.14 $\pm$ 0.19} & 3 \\
\bottomrule
\end{tabular}
\end{table*}

\begin{figure}[t]\centering\includegraphics[width=\columnwidth]{bars_main_miou.pdf}
\caption{mIoU of the four systems at $N=500$ (mean and spread over seeds; generated from the \texttt{main} group).}\label{fig:bars}\end{figure}

\begin{figure}[t]\centering\includegraphics[width=0.95\columnwidth]{size_curves.pdf}
\caption{CE test mIoU against training-set size (mean $\pm$ std over seeds). Boundary noise was run only at $N=100,500,2000$.}\label{fig:size}\end{figure}

\begin{figure}[t]\centering\includegraphics[width=0.95\columnwidth]{perclass_ce.pdf}
\caption{Per-class test IoU of CE at $N=500$ for clean, boundary-noise and flip-noise training labels.}\label{fig:perclass}\end{figure}

We go hypothesis by hypothesis. All tests are the Welch $t$-tests registered in \texttt{proposal.md}, with four seeds per arm and no multiplicity correction; $p$-values below are descriptive given the very small $n$.

\textbf{H1 (more clean data helps): supported.} CE on clean labels reaches mIoU $\rhval{analysis/analysis/flip_p0.3/ce_miou_none_n100/mean:2}$ ($\pm\rhval{analysis/analysis/flip_p0.3/ce_miou_sd_none_n100/mean:2}$) at $N=100$ and $\rhval{analysis/analysis/flip_p0.3/ce_miou_none_n2000/mean:2}$ ($\pm\rhval{analysis/analysis/flip_p0.3/ce_miou_sd_none_n2000/mean:2}$) at $N=2000$; the difference is $\rhval{analysis/analysis/flip_p0.3/h1_diff/mean:2}$ (Welch $p=\rhval{analysis/analysis/flip_p0.3/h1_p/mean:3}$). The curve (Fig.~\ref{fig:size}, Table~\ref{tab:size}) is concave: most of the gain is reached by $N=500$ (\rhval{analysis/analysis/flip_p0.3/ce_miou_none_n500/mean:2}), and $N=1000$ gives \rhval{analysis/analysis/flip_p0.3/ce_miou_none_n1000/mean:2}. Square and triangle IoU rise most (Table~\ref{tab:size}).

\textbf{H2 (flips hurt more than boundary jitter at $N=500$): supported.} CE obtains \rhval{analysis/analysis/flip_p0.3/ce_miou_boundary_p03_n500/mean:2} under boundary noise and \rhval{analysis/analysis/flip_p0.3/ce_miou_flip_p03_n500/mean:2} under flip noise at the same nominal level; the difference is \rhval{analysis/analysis/flip_p0.3/h2_diff/mean:2} (Welch $p=\rhval{analysis/analysis/flip_p0.3/h2_p/mean:3}$). Boundary noise is not distinguishable from clean labels at $N=500$ (drop \rhval{analysis/analysis/flip_p0.3/drop_boundary_n500/mean:2}, Welch $p=\rhval{analysis/analysis/flip_p0.3/h2_boundary_vs_clean_p/mean:2}$). The comparison is confounded by the fact that ``the same nominal level'' touches very different numbers of pixels (the fraction of training pixels whose label is changed, \texttt{label\_noise\_pixels}, is \rhval{main/ce/boundary_p0.3/label_noise_pixels/mean:3} for boundary noise and \rhval{main/ce/flip_p0.3/label_noise_pixels/mean:3} for flips at $N=500$), so the result says the two \emph{models at $p=0.3$} differ, not that a flipped pixel is intrinsically more harmful than a jittered one. Fig.~\ref{fig:perclass} shows that flips damage square and triangle most (CE square IoU \rhval{analysis/analysis/flip_p0.3/ce_iou_square_none_n500/mean:2} clean vs \rhval{analysis/analysis/flip_p0.3/ce_iou_square_flip_p03_n500/mean:2} under flips; triangle \rhval{analysis/analysis/flip_p0.3/ce_iou_triangle_none_n500/mean:2} vs \rhval{analysis/analysis/flip_p0.3/ce_iou_triangle_flip_p03_n500/mean:2}), circle less (\rhval{analysis/analysis/flip_p0.3/ce_iou_circle_none_n500/mean:2} vs \rhval{analysis/analysis/flip_p0.3/ce_iou_circle_flip_p03_n500/mean:2}), and leave background IoU essentially unchanged (\rhval{analysis/analysis/flip_p0.3/ce_iou_bg_none_n500/mean:2} vs \rhval{analysis/analysis/flip_p0.3/ce_iou_bg_flip_p03_n500/mean:2}); the reported mIoU therefore hides that the damage is entirely a foreground-class confusion.

\textbf{H3 (flip damage is larger with less data): refuted / inconclusive, opposite direction.} The drop caused by flips (clean minus flip mIoU, mean over four seeds) is \rhval{analysis/analysis/flip_p0.3/drop_flip_p03_n100/mean:2} at $N=100$ and \rhval{analysis/analysis/flip_p0.3/drop_flip_p03_n2000/mean:2} at $N=2000$; the difference of drops, $N=100$ minus $N=2000$, is \rhval{analysis/analysis/flip_p0.3/h3_flip_p03_diff/mean:2} (Welch $p=\rhval{analysis/analysis/flip_p0.3/h3_flip_p03_p/mean:2}$). The registered direction is not supported: there is no evidence that more data dilutes object-flip noise, and the point estimate goes the other way. Over seeds 0--2 shared by all sizes the drop is \rhval{analysis/analysis/flip_p0.3/drop_flip_s012_n100/mean:2}, \rhval{analysis/analysis/flip_p0.3/drop_flip_s012_n250/mean:2}, \rhval{analysis/analysis/flip_p0.3/drop_flip_s012_n500/mean:2}, \rhval{analysis/analysis/flip_p0.3/drop_flip_s012_n1000/mean:2}, \rhval{analysis/analysis/flip_p0.3/drop_flip_s012_n2000/mean:2} for $N=100,250,500,1000,2000$, increasing at first and then flat. For boundary noise the drops are tiny at both sizes (\rhval{analysis/analysis/flip_p0.3/drop_boundary_p03_n100/mean:3} and \rhval{analysis/analysis/flip_p0.3/drop_boundary_p03_n2000/mean:3}; difference test $p=\rhval{analysis/analysis/flip_p0.3/h3_boundary_p03_p/mean:2}$). With flips, the seed-to-seed std of CE grows with $N$ (\rhval{analysis/analysis/flip_p0.3/ce_miou_sd_flip_p03_n100/mean:2} at $N=100$, \rhval{analysis/analysis/flip_p0.3/ce_miou_sd_flip_p03_n2000/mean:2} at $N=2000$): some seeds recover a decent triangle class and others do not, and four seeds cannot resolve this.

\textbf{H4 (robust losses beat CE under noise and do not lose on clean data): not supported.} Table~\ref{tab:cmp} lists the registered comparisons (\texttt{rh compare}, CE as reference). GCE equals CE on clean and boundary-noise labels and has a higher mean under flips, but the effect is not significant (CE minus GCE \rhval{cmp/main/gce/flip_p0.3/miou/delta:2}, $p=\rhval{cmp/main/gce/flip_p0.3/miou/welch_p:2}$); a seed-paired test, which we do not use for the verdict, gives $p=\rhval{cmp/main/gce/flip_p0.3/miou/paired_p:3}$, so the direction is consistent across seeds but its size is not established. SCE is not better than CE on any task and is worse on boundary noise in mean (CE minus SCE \rhval{cmp/main/sce/boundary_p0.3/miou/delta:2}, $p=\rhval{cmp/main/sce/boundary_p0.3/miou/welch_p:2}$), where its triangle IoU is \rhval{main/sce/boundary_p0.3/iou_triangle/mean:2} with a std of \rhval{main/sce/boundary_p0.3/iou_triangle/std:2} over seeds. Our band-ignore CE is clearly worse than CE on clean labels (difference \rhval{cmp/main/band-ignore-ce/none/miou/delta:2}, $p=\rhval{cmp/main/band-ignore-ce/none/miou/welch_p:3}$) and on boundary noise (\rhval{cmp/main/band-ignore-ce/boundary_p0.3/miou/delta:2}, $p=\rhval{cmp/main/band-ignore-ce/boundary_p0.3/miou/welch_p:3}$), i.e. the method designed for boundary noise fails exactly there; under flips it is lower in mean (\rhval{cmp/main/band-ignore-ce/flip_p0.3/miou/delta:2}, $p=\rhval{cmp/main/band-ignore-ce/flip_p0.3/miou/welch_p:2}$). A plausible but untested reason is that, in a $64\times64$ image with small objects, a two-pixel band removes supervision on a large share of the foreground pixels and on the true edges that define square and triangle corners; we did not measure the ignored fraction.

\begin{table}[t]\centering\caption{Registered comparisons against CE at $N=500$ (mIoU; difference is CE minus system, so negative means the system is better; four seeds, Welch $p$).}\label{tab:cmp}
\begin{tabular}{llrr}\toprule Task & System & CE$-$system & Welch $p$\\\midrule
clean & GCE & \rhval{cmp/main/gce/none/miou/delta:3} & \rhval{cmp/main/gce/none/miou/welch_p:3}\\
clean & SCE & \rhval{cmp/main/sce/none/miou/delta:3} & \rhval{cmp/main/sce/none/miou/welch_p:3}\\
clean & Band-ignore CE & \rhval{cmp/main/band-ignore-ce/none/miou/delta:3} & \rhval{cmp/main/band-ignore-ce/none/miou/welch_p:3}\\
boundary 0.3 & GCE & \rhval{cmp/main/gce/boundary_p0.3/miou/delta:3} & \rhval{cmp/main/gce/boundary_p0.3/miou/welch_p:3}\\
boundary 0.3 & SCE & \rhval{cmp/main/sce/boundary_p0.3/miou/delta:3} & \rhval{cmp/main/sce/boundary_p0.3/miou/welch_p:3}\\
boundary 0.3 & Band-ignore CE & \rhval{cmp/main/band-ignore-ce/boundary_p0.3/miou/delta:3} & \rhval{cmp/main/band-ignore-ce/boundary_p0.3/miou/welch_p:3}\\
flip 0.3 & GCE & \rhval{cmp/main/gce/flip_p0.3/miou/delta:3} & \rhval{cmp/main/gce/flip_p0.3/miou/welch_p:3}\\
flip 0.3 & SCE & \rhval{cmp/main/sce/flip_p0.3/miou/delta:3} & \rhval{cmp/main/sce/flip_p0.3/miou/welch_p:3}\\
flip 0.3 & Band-ignore CE & \rhval{cmp/main/band-ignore-ce/flip_p0.3/miou/delta:3} & \rhval{cmp/main/band-ignore-ce/flip_p0.3/miou/welch_p:3}\\
\bottomrule\end{tabular}
\end{table}

\textbf{Failure cases.} (i) Band-ignore CE as above. (ii) Triangle IoU of GCE and CE under flips has a per-seed spread of several tenths (Table~\ref{tab:main}), because in some seeds the triangle class is learned and in others it is largely predicted as another class. (iii) SCE has the lowest clean-label mean of the three published-loss systems (CE minus SCE \rhval{cmp/main/sce/none/miou/delta:2}, $p=\rhval{cmp/main/sce/none/miou/welch_p:2}$); the reverse-CE term may make optimisation slower in our short budget, so the baseline may be disadvantaged, but we did not test this and did not retune its defaults.
